% Physical unit 010; accepted Session 02.
\GordanUnitStart{010}{printed page 5}
\GordanTopFolio{5}
\[
 \delta\int\left\{\sqrt{dx^2+dy^2+dz^2}+\lambda f\right\}=0,
\]
or, if instead of the rectangular coordinates \gq{$x,y,z$} we use the elliptic coordinates $l,\psi$,
\[
 \delta\int \sqrt{(1-e^2\cos l)\,dl^2+\cos^2l\,d\psi^2}=0.
\]

\begin{center}\S\ 4.\end{center}
Here, if we apply the customary methods of analysts, one finally obtains
\[
 {\cos^2l\,{d\psi\over dl}\over\sqrt{1-e^2\cos^2l+\cos^2l\,{d\psi\over dl}{d\psi\over dl}}}=\mathrm{Const}.
\]
Using in this formula
\[
 {\cos l\,{d\psi\over dl}\over\sqrt{1-e^2\cos^2l+\cos^2l\,{d\psi\over dl}{d\psi\over dl}}}=\sin T
\]
\hfill compare Chapter I, \S\ 9,

we have
\[
 \cos l\sin T=\mathrm{Const}.
\]
And if, taking
\[
 T={\pi\over 2},\qquad l=l_0,
\]
we put
\[
 \cos l\sin T=\cos l_0,
\]
then
\[
 \cos l\cos T=\sqrt{\sin^2l_0-\sin^2l},
\qquad
 \tg T={\cos l_0\over\sqrt{\sin^2l_0-\sin^2l}}.
\]

\begin{center}\S\ 5.\end{center}
That is: a geodetic line which touches some circle also touches another circle having the corresponding position. Further, since always
\[
 l_0>l,
\]
it can never happen that the geodetic line recedes from the equator beyond these circles. Therefore we call the geodetic line tangent to the circle $l_0$ the geodetic line $l_0$, and the point at which the circle is touched the origin of the geodetic line; from this origin we shall henceforth measure the extensions \gq{$s$} and the longitudes $\psi$ of the points.

\begin{center}\S\ 6.\end{center}
To determine these from the latitude $l$, applying formulas II to the formulas proposed in Chapter I, \S\ 9, we have
